\section{Architettura del codice}\label{architettura-del-codice}

A livello software, BlueAgent organizza l\textquotesingle interazione
con l\textquotesingle host disaccoppiando l\textquotesingle API in
\passthrough{\lstinline!/api!} dalla gestione a basso livello dei
container e delle componenti di rete, gestita ad alto livello dai
servizi in \passthrough{\lstinline!/services!}, che contiene la logica
di business, mentre le complessità operative sono delegate a moduli e
driver specializzati (\textbf{aggiungi cartelle}).

\begin{center}\rule{0.5\linewidth}{0.5pt}\end{center}

\subsection{\texorpdfstring{Astrazione del Runtime:
\texttt{ContainerEngine}}{Astrazione del Runtime: ContainerEngine}}\label{astrazione-del-runtime-containerengine}

BlueAgent non invoca direttamente il demone Incus
all\textquotesingle interno della logica delle challenge. Tutte le
operazioni di gestione del ciclo di vita dei container sono intermediate
da un\textquotesingle interfaccia astratta denominata
\passthrough{\lstinline!ContainerEngine!}.

Questa classe astratta formalizza esclusivamente le primitive necessarie
ai requisiti operativi del Cyber Range:

\begin{itemize}
\tightlist
\item
  Creazione, avvio, arresto forzato ed eliminazione delle istanze.
\item
  Duplicazione rapida e gestione delle interfacce di rete.
\item
  Esecuzione dei comandi nei container.
\end{itemize}

Disaccoppiare l\textquotesingle interfaccia
dall\textquotesingle implementazione permette ai servizi di non dover
conoscere come avviene la gestione dei container, e di predisporre il
codice per poter utilizzare altri gestori.

\begin{quote}
\textbf{ATTENZIONE: anche la generazione delle immagini dovrebbe
avvenire tramite interfaccia}
\end{quote}

L\textquotesingle implementazione concreta
\passthrough{\lstinline!IncusEngine!} comunica con il demone di sistema
tramite Unix Domain Socket
(\passthrough{\lstinline!/var/lib/incus/unix.socket!}), traducendo le
primitive astratte in chiamate REST HTTP asincrone.

\begin{quote}
\passthrough{\lstinline!Unix Domain Socket (UDS)!}: Meccanismo di
Inter-Process Communication (IPC) gestito interamente dal kernel del
sistema operativo. La comunicazione avviene trasferendo i dati
direttamente tra i buffer di memoria dei processi nello spazio kernel
senza attraversare lo stack di rete. L\textquotesingle endpoint è
rappresentato da un file speciale sul filesystem (es:
\passthrough{\lstinline!/var/lib/incus/unix.socket!}).{[}\^{}7{]}
\end{quote}

\begin{center}\rule{0.5\linewidth}{0.5pt}\end{center}

\subsection{\texorpdfstring{Ciclo di Vita della Sandbox: Async Context
Manager
(\texttt{OVNVerifier})}{Ciclo di Vita della Sandbox: Async Context Manager (OVNVerifier)}}\label{ciclo-di-vita-della-sandbox-async-context-manager-ovnverifier}

Il processo di verifica richiede di istanziare una replica fedele della
rete di produzione che deve esistere esclusivamente per la durata dei
test. In caso di qualsiasi tipo di fallimento,
l\textquotesingle infrastruttura clonata deve essere rilasciata per
evitare risorse allocate e sprecate.

Per garantire questa proprietà, il ciclo di vita della sandbox è
incapsulato all\textquotesingle interno della classe
\passthrough{\lstinline!OVNVerifier!}, implementata come
\textbf{Asynchronous Context Manager} di Python.

\begin{quote}
\textbf{Pattern RAII (Resource Acquisition Is Initialization):}
L\textquotesingle allocazione delle risorse è vincolato
dall\textquotesingle acquisizione del contesto, mentre il loro rilascio
e smantellamento deterministico è garantito all\textquotesingle uscita
dal blocco, gestendo automaticamente qualsiasi anomalia o interruzione.
{[}\^{}8{]}
\end{quote}

Il ciclo di vita dell\textquotesingle{}
\passthrough{\lstinline!OVNVerifier!} segue tre fasi:

\begin{enumerate}
\def\labelenumi{\arabic{enumi}.}
\tightlist
\item
  \textbf{Inizializzazione (\passthrough{\lstinline!\_\_aenter\_\_!}):}

  \begin{itemize}
  \tightlist
  \item
    Creazione della rete OVN temporanea
  \item
    Duplicazione dei container tramite snapshot ZFS
  \item
    Riconnessione dell\textquotesingle interfaccia
    \passthrough{\lstinline!eth0!} dei cloni alla nuova rete OVN
  \item
    Istanziazione opzionale del container \emph{Verifier Runner}.
  \end{itemize}
\item
  \textbf{Esecuzione delle Verifiche
  (\passthrough{\lstinline!run\_checks!}):}

  \begin{itemize}
  \tightlist
  \item
    Esecuzione controllata dei comandi di test con timeout
  \end{itemize}
\item
  \textbf{Rilascio e Teardown
  (\passthrough{\lstinline!\_\_aexit\_\_!}):}

  \begin{itemize}
  \tightlist
  \item
    Smantellamento delle risorse:

    \begin{enumerate}
    \def\labelenumii{\arabic{enumii}.}
    \tightlist
    \item
      Arresto ed eliminazione del container
      \passthrough{\lstinline!verify-runner!};
    \item
      Arresto forzato ed eliminazione di tutti i container cloni;
    \item
      Eliminazione della rete virtuale OVN.
    \end{enumerate}
  \end{itemize}
\end{enumerate}

L\textquotesingle uso dell\textquotesingle{}\textbf{Async Context
Manager} fornisce dunque un metodo sicuro per utilizzare
l\textquotesingle oggetto \passthrough{\lstinline!OVNVerifier!}

\begin{center}\rule{0.5\linewidth}{0.5pt}\end{center}

\subsection{Concorrenza e Two-Tier
Locking}\label{concorrenza-e-two-tier-locking}

Incus non supporta operazioni di mutazione concorrenti sul medesimo
container: qualora giungano due richieste simultanee, il demone Incus
rifiuta la transazione. Inoltre potrebbero nascere altri problemi
nell\textquotesingle accesso alle risorse dei team

Per prevenire questo fenomeno, BlueAgent implementa un meccanismo di
sincronizzazione a due livelli (\passthrough{\lstinline!TwoTierLock!}):

\begin{itemize}
\tightlist
\item
  \textbf{Tier 1 (Intra-processo):} utilizza
  \passthrough{\lstinline!asyncio.Lock!} per coordinare e serializzare
  le coroutine concorrenti all\textquotesingle interno del singolo event
  loop.
\item
  \textbf{Tier 2 (Inter-processo):} applica un lock non bloccante sul
  filesystem (\passthrough{\lstinline!fcntl.flock!}) memorizzato nel
  file \passthrough{\lstinline!/tmp/blueagent\_locks/<risorsa>.lock!}.
\end{itemize}

Questa architettura protegge le risorse sia a livello di singolo
processo sia in presenza di server ASGI multi-worker (come Uvicorn o
Gunicorn configurati con molteplici worker di sistema), garantendo mutua
esclusione sulle mutazioni dei container senza bloccare
l\textquotesingle event loop asincrono.

\begin{center}\rule{0.5\linewidth}{0.5pt}\end{center}

\subsection{Pipeline Asincrona non Bloccante su Socket
Unix}\label{pipeline-asincrona-non-bloccante-su-socket-unix}

Le operazioni di virtualizzazione (clonazione di storage, avvio di init
system complessi nei container di sistema, polling delle operazioni)
comportano latenze nell\textquotesingle ordine di decine o centinaia di
millisecondi. L\textquotesingle adozione di un modello sincrono
multithread comporterebbe un consumo eccessivo di risorse e colli di
bottiglia sotto carico.

BlueAgent è progettato secondo un paradigma completamente asincrono:

\begin{itemize}
\tightlist
\item
  \passthrough{\lstinline!IncusEngine!} comunica con il socket UDS
  locale (\passthrough{\lstinline!/var/lib/incus/unix.socket!}) per
  mezzo di \passthrough{\lstinline!httpx.AsyncClient!}.
\item
  Poiché le operazioni di mutazione in Incus sono gestite in background
  dal demone e restituiscono un UUID di tracciamento, BlueAgent
  interroga asincronamente l\textquotesingle endpoint
  \passthrough{\lstinline!/1.0/operations/<uuid>/wait!}, evitando cicli
  di polling bloccanti.
\item
  Le chiamate verso il kernel Linux (\passthrough{\lstinline!nftables!},
  \passthrough{\lstinline!ip link!}, WireGuard) sfruttano
  \passthrough{\lstinline!asyncio.create\_subprocess\_exec!}.
\end{itemize}

Questo approccio consente a un singolo processo BlueAgent di servire le
richieste di telemetria e validazione di centinaia di utenti con un
overhead di memoria e CPU quasi nullo.

\begin{center}\rule{0.5\linewidth}{0.5pt}\end{center}

\subsection{\texorpdfstring{Resilienza e Riconciliazione dello Stato
(\texttt{LifecycleManager})}{Resilienza e Riconciliazione dello Stato (LifecycleManager)}}\label{resilienza-e-riconciliazione-dello-stato-lifecyclemanager}

Qualora l\textquotesingle host subisca un riavvio forzato improvviso o
il processo di BlueAgent venga terminato bruscamente con
\passthrough{\lstinline!SIGKILL!}, le routine di pulizia
(\passthrough{\lstinline!\_\_aexit\_\_!}) non possono intervenire,
lasciando container e switch virtuali orfani sul sistema.

Per risolvere queste anomalie senza intervento manuale, BlueAgent
include il demone in background
\passthrough{\lstinline!LifecycleManager!}:

\begin{itemize}
\tightlist
\item
  Scansiona periodicamente le istanze e le reti con prefisso
  \passthrough{\lstinline!verify-*!} presenti su Incus, confrontandole
  con lo stato persistito nel database e forzandone
  l\textquotesingle eliminazione se non risultano associate a sessioni
  attive.
\item
  \textbf{Grace Period di 90 Secondi:} il demone ignora qualsiasi
  risorsa la cui creazione risalga a meno di 90 secondi. Questa finestra
  temporale risolve una race condition fondamentale: durante
  l\textquotesingle avvio di una verifica, la rete OVN viene creata
  prima che i container cloni vi siano collegati. Senza il periodo di
  grazia, il ciclo di bonifica interpreterebbe la rete appena creata
  come una risorsa "vuota e orfana", distruggendola mentre la richiesta
  legittima ne sta ancora completando la configurazione.
\end{itemize}

\begin{center}\rule{0.5\linewidth}{0.5pt}\end{center}
